\section*{第 \ref{ch:b2:muscle-movement} 章——肌肉收缩与运动功能}

\begin{solution}{exo:b2:muscle-movement:1}
细丝沿粗丝滑动，由\omterm{prop:b2:muscle-movement:crossbridge}{肌球蛋白}头部牵引；
两种丝都不缩短。I 带与 H 区变窄，Z 盘
相互靠近；A 带（即粗丝）长度保持不变。
\end{solution}

\begin{solution}{exo:b2:muscle-movement:2}
（1）ATP 结合，头部脱离；ATP 水解，头部扳起。（2）头部
与肌动蛋白结合。（3）磷酸释放，发生动力冲程，ADP 释放。
（4）新的 ATP 结合，头部脱离。ATP 是使头部
\emph{脱离}所必需的：若无 ATP，所有头部都会保持结合，
肌肉锁死，即僵直。
\end{solution}

\begin{solution}{exo:b2:muscle-movement:3}
神经冲动；乙酰胆碱释放并产生终板电位
（\qty{0.6}{ms}）；肌肉\omterm{prop:b2:neurons-synapses:ap}{动作电位}沿\omterm{def:b2:cell-differentiation:myogenesis}{肌纤维}传播，并沿
T 小管下行（\qty{1}{ms}）；\omterm{def:b2:cell-differentiation:fibre}{肌质网}释放钙离子
（\qty{1}{ms}）；钙离子与\omterm{prop:b2:muscle-movement:coupling}{肌钙蛋白}结合，\omterm{prop:b2:muscle-movement:coupling}{原肌球蛋白}移动，
头部附着（数毫秒）；力在数十
毫秒内上升。
\end{solution}

\begin{solution}{exo:b2:muscle-movement:4}
\omterm{def:b2:muscle-movement:motorunit}{运动单位}：一个运动\omterm{def:b2:neurons-synapses:neuron}{神经元}及其支配的全部\omterm{def:b2:cell-differentiation:myogenesis}{肌纤维}。
\omterm{def:b2:muscle-movement:motorunit}{单收缩}：单次冲动引起的短暂收缩。\omterm{def:b2:muscle-movement:motorunit}{强直收缩}：冲动
到达速度快于\omterm{def:b2:muscle-movement:motorunit}{单收缩}衰减速度时形成的融合持续收缩。
力的分级取决于每个单位的放电频率，
以及被\omterm{def:b2:muscle-movement:motorunit}{募集}的单位数目。
\end{solution}

\begin{solution}{exo:b2:muscle-movement:5}
每个半\omterm{def:b2:cell-differentiation:fibre}{肌节}在 \qty{50}{ms} 内滑动 \qty{0.2}{\micro m}：
\qty{4}{\micro m/s}；相当于 20 次 \qty{10}{nm} 的冲程，
在 \qty{50}{ms} 内完成，若某一头部始终附着，则每秒 400 次。
\end{solution}

\begin{solution}{exo:b2:muscle-movement:6}
\qty{2}{cm} 用时 \qty{0.1}{s}：\qty{0.2}{m/s}，即每秒 2 个长度；每个
\omterm{def:b2:cell-differentiation:fibre}{肌节}缩短 $2/\num{40000}$ cm $= \qty{0.5}{\micro m}$，在
\qty{0.1}{s} 内，即 \qty{5}{\micro m/s}。
\end{solution}

\begin{solution}{exo:b2:muscle-movement:7}
$F/F_0 = 0.3125/(v/v_{\max} + 0.25) - 0.25$：速度为每秒 1、3、6 个
长度时（即 $0.1$、$0.3$、$0.6$ 倍的 $v_{\max}$），相对力分别为 $0.64$、$0.32$、$0.12$。
功率 $Fv$（以 $F_0\times$ 每秒长度数为单位）：$0.64$、$0.95$、$0.71$
——在每秒约 3 个长度、约 $v_{\max}$ 三分之一处取最大值。
\end{solution}

\begin{solution}{exo:b2:muscle-movement:8}
$80\times 30 = \qty{2400}{N}$；在脚部为 $2400\times 5/45 =
\qty{270}{N}$。
\end{solution}

\begin{solution}{exo:b2:muscle-movement:9}
在 \qty{10}{Hz} 时，冲动每 \qty{100}{ms} 到达一次，此时上一次的
钙离子已消失：形成各自独立的\omterm{def:b2:muscle-movement:motorunit}{单收缩}，每次都从部分松弛的状态开始，
表现为波纹状起伏。在 \qty{50}{Hz} 时，冲动每
\qty{20}{ms} 到达一次，钙离子尚未被泵走，
钙离子浓度保持较高，力发生融合。\omterm{def:b2:muscle-movement:motorunit}{强直收缩}的力
大于\omterm{def:b2:muscle-movement:motorunit}{单收缩}，是因为单次钙离子脉冲过于短暂，
横桥不足以拉伸串联弹性元件并产生完全的力；
持续的钙离子则可使其做到这一点。
\end{solution}

\begin{solution}{exo:b2:muscle-movement:10}
做功 \qty{20}{kJ}，ATP 所含能量 \qty{80}{kJ}，即 $1.6$ mol 的 ATP。最初
\qty{5}{s} 内：$0.8$ mol \omterm{prop:b2:muscle-movement:energy}{磷酸肌酸}。其余 $0.8$ mol 的 ATP
由糖酵解提供：$0.4$ mol 葡萄糖，生成 $0.8$ mol 乳酸——
在 \qty{40}{L} 中，浓度为 \qty{20}{mmol/L}。
\end{solution}

\begin{solution}{exo:b2:muscle-movement:11}
释放通道被阻断：无钙离子，无收缩，即麻痹。
泵被阻断：钙离子持续升高，肌肉无法舒张，即
挛缩。\omterm{prop:b2:muscle-movement:coupling}{肌钙蛋白}不敏感：\omterm{prop:b2:muscle-movement:coupling}{原肌球蛋白}永不移动，无
收缩。\omterm{def:b2:blood-circulation:blood}{血液}中无钙离子：骨骼肌仍可正常收缩
（其钙离子来自细胞内）；\omterm{def:b2:heart:anatomy}{心脏}则需要
细胞外钙离子触发其释放，因而停止跳动。
\end{solution}

\begin{solution}{exo:b2:muscle-movement:12}
头部最多可将 ATP 自由能的一半转化为功；
其余为热量，此外还有泵的损耗、以较快速度附着却不产生拉力
的横桥的损耗，以及弹性元件的损耗，
使总效率仅约四分之一——与优质汽油机的效率相当。
两种分级力的方式兼顾两方面优点：低强度时可精细调控
（小单位逐个增加，各自放电频率加快或减慢）；
同时力的范围很大（约千倍，从一个小单位在\omterm{def:b2:muscle-movement:motorunit}{单收缩}频率下工作，
到所有单位处于\omterm{def:b2:muscle-movement:motorunit}{强直收缩}）。若仅有一个开关，则只能得到
仅一种力或完全没有力。
\end{solution}

\begin{solution}{pb:b2:muscle-movement:1}
\textbf{1.} $v = \sqrt{2\times 9.81\times 0.4} = \qty{2.8}{m/s}$；
$\tfrac12\times 70\times 2.8^{2} = \qty{275}{J}$。
\textbf{2.} 做功 $= 275 + 70\times 9.81\times 0.4 = \qty{550}{J}$，
作用距离 \qty{0.4}{m}：\qty{1370}{N}，为体重的两倍。
\textbf{3.} $F_0 = 500\times 30 = \qty{15000}{N}$；在脚部为
\qty{1500}{N}。
\textbf{4.} \qty{4}{cm} 用时 \qty{0.3}{s}：\qty{0.13}{m/s}，即每秒 $1.1$ 个
长度，$0.14\,v_{\max}$。
\textbf{5.} $F/F_0 = 0.3125/(0.14 + 0.25) - 0.25 = 0.55$：肌肉产生 \qty{8300}{N}，
在脚部为 \qty{830}{N}——远低于所需的
\qty{1370}{N}。仅靠肌肉缩短无法实现这一跳跃。
预先下蹲的动作提供其余部分：跳跃者下蹲时，
被拉伸的肌腱储存弹性能；肌肉在高力下几乎等长收缩，
使肌腱承受负荷；随后肌腱回弹，
速度快于任何\omterm{def:b2:cell-differentiation:myogenesis}{肌纤维}所能达到的缩短速度。
\textbf{6.} $550/0.3 = \qty{1.8}{kW}$；伸肌每千克 \qty{92}{W/kg}。
\textbf{7.} $4/\num{48000}$ cm $= \qty{0.83}{\micro m}$，每个\omterm{def:b2:cell-differentiation:fibre}{肌节}；
每半\omterm{def:b2:cell-differentiation:fibre}{肌节} \qty{0.42}{\micro m}：共 42 个 \qty{10}{nm} 的冲程。
\textbf{8.} \qty{0.42}{\micro m} 用时 \qty{0.3}{s}：\qty{1.4}{\micro m/s}；
若持续附着，则每秒 140 个冲程。
\textbf{9.} $500\times 6\times 10^{10}\times \num{48000}\times 300 =
4.3\times 10^{20}$ 个头部。
\textbf{10.} $0.55\times \num{15000}\times 0.04 = \qty{330}{J}$；每个
冲程 $0.4\times 8\times 10^{-20} = \qty{3.2e-20}{J}$；$1.0\times
10^{22}$ 个冲程。
\textbf{11.} $1.0\times 10^{22}/4.3\times 10^{20} = 24$ 个冲程
每个头部，占可用 42 个的约 \qty{60}{\%}。余量
很小——肌肉正接近其工作极限，
正如第 5 题所示。
\textbf{12.} $1.0\times 10^{22}/6.0\times 10^{23} = 0.017$ mol，
\qty{8.7}{g}；其能量 $0.017\times 50 = \qty{860}{J}$，其中
\qty{330}{J} 转化为功：\qty{39}{\%}。
\textbf{13.} 冲动频率需高于每 \qty{30}{ms} 一次，即高于
\qty{33}{Hz}；\qty{0.3}{s} 内约 10 个冲动。
\textbf{14.} $9.9\times 10^{-6}\times 5\times 10^{-10}\times 6\times
10^{23} = 3\times 10^{9}$ 个游离离子；$1.5\times 10^{9}$ 个 ATP。
\textbf{15.} 纤维数：$500/(\pi\times(3\times 10^{-3})^{2}) = 1.8\times
10^{7}$；每根纤维每次冲动的横桥 ATP 消耗 $1.0\times
10^{22}/(1.8\times 10^{7}\times 10) = 6\times 10^{13}$，为泵耗估算值的四万倍。
游离钙离子只是很小的一部分：\omterm{def:b2:cell-differentiation:fibre}{肌质网}
释放的钙离子量约为 \qty{0.1}{mmol/L} 的数量级，
是游离钙离子升高量的十倍，几乎全部被\omterm{prop:b2:muscle-movement:coupling}{肌钙蛋白}和缓冲液同时结合——
而这些钙离子都必须被泵回。
\textbf{16.} 对于最大跳跃，几乎所有单位都须在
数十毫秒内被\omterm{def:b2:muscle-movement:motorunit}{募集}，先小后大；轻微的一步仅\omterm{def:b2:muscle-movement:motorunit}{募集}其中约
十分之一至五分之一。
\textbf{17.} \omterm{def:b2:muscle-movement:motorunit}{运动单位}按大小顺序被\omterm{def:b2:muscle-movement:motorunit}{募集}：低强度时，每新增一个单位
只带来很小的增量；高强度时，剩余单位较大，
每个都带来较大的一步。
\textbf{18.} 每根纤维的终板电位均减半；
电位降至阈值以下的纤维不再放电，因此
\omterm{def:b2:muscle-movement:motorunit}{单收缩}与\omterm{def:b2:muscle-movement:motorunit}{强直收缩}均变小，跳跃高度不足。
\textbf{19.} $5\times 20 = \qty{0.1}{mol}$ 的 ATP：相当于六次跳跃的用量。
\textbf{20.} $20\times 20 = \qty{0.4}{mol}$ 的\omterm{prop:b2:muscle-movement:energy}{磷酸肌酸}：
约可供 23 次跳跃。
\textbf{21.} 十次跳跃消耗 $0.17$ mol，尚未耗尽
\omterm{prop:b2:muscle-movement:energy}{磷酸肌酸}；约二十次之后，糖酵解接替供能，
乳酸与质子随之积累。
\textbf{22.} \qty{1.8}{kW} 的做功，效率为 \qty{40}{\%}，即需 \qty{4.6}{kW} 的 ATP 能量，
即每秒 $0.09$ mol，每秒需氧气 $0.015$ mol——
\qty{0.34}{L/s}，\qty{21}{L/min}，是\omterm{def:b2:blood-circulation:blood}{血液}所能输送量的十倍：
跳跃在本质上是无氧的，氧气需事后补充。
\textbf{23.} $0.17$ mol 的 ATP 需要 $0.029$ mol 的氧气，
\qty{0.64}{L}；若在一分钟内偿还，则为 \qty{640}{mL/min}，是
静息耗氧量的两倍半。
\textbf{24.} ATP 可在数毫秒内由\omterm{prop:b2:muscle-movement:energy}{磷酸肌酸}再磷酸化，
因此其浓度几乎不下降；僵直需要 ATP
接近于零，而任何仍有储备的活体肌肉都不会达到这一水平。
\textbf{25.} 每个半\omterm{def:b2:cell-differentiation:fibre}{肌节}可用 42 个冲程，需要 24 个；
跳跃速度下 $F/F_0 = 0.55$；ATP 为 $0.017$ mol，\qty{8.7}{g}。
\end{solution}
